% preamble_build.tex: copy of Maximus's preamble.tex for the DE v3 build.
% The blocks marked [BUILD COPY CHANGE] are the earlier build fixes.
% The blocks marked [V3 ADDITION] are new in this version (example box,
% spacing). Everything else is unchanged from preamble.tex.

% For margins
\usepackage[a4paper, margin = 20mm]{geometry}

\usepackage[Lenny]{fncychap}
% Rejne, Glenn, Lenny

\usepackage{etoc}

% [V3 ADDITION] cancel: used to show the dx's "cancelling" in Section 11.3
\usepackage{cancel}
\renewcommand{\CancelColor}{\color{red}}

% For enumeration & math symbols
\usepackage{enumerate}
\usepackage{amsmath,amsfonts,amssymb,amsthm,graphicx,lipsum}

% For diagrams, plots and boxes
\usepackage{pgfplots}
\pgfplotsset{compat=1.18}
\usepackage{tikz}
\usepackage[most]{tcolorbox}
\tcbuselibrary{skins}
\usepackage{xcolor}
\pgfplotsset{
    standard/.style={clip = false,
    axis line style = thick,
    trig format=rad,
    axis lines = middle,
    every axis x label/.style={at={(current axis.right of origin)},anchor= west},
    every axis y label/.style={at={(current axis.above origin)},anchor=south}
    }
}

% For GC font
\usepackage{fontspec}
% [BUILD COPY CHANGE 1] The TI-84 key font is loaded only if the file exists
% (checked in ./Fonts/... and ../Fonts/...). Otherwise \ti{} falls back to a
% small boxed sans-serif key. Known letters of the TI font are mapped to their
% key names in the fallback: o = [y=], z = [mode], f = [L3].
\newcommand{\tikey}[1]{\tcbox[on line, enhanced, boxsep=0pt, left=2pt, right=2pt,
  top=1pt, bottom=1pt, boxrule=0.5pt, arc=1.5pt, colback=white,
  colframe=black!70]{\sffamily\footnotesize\strut #1}}
\ExplSyntaxOn
\NewDocumentCommand{\tifallback}{m}{
  \str_case:nnF {#1}
    { {o}{\tikey{y=}} {z}{\tikey{mode}} {f}{\tikey{L3}} }
    { \tikey{#1} }
}
\ExplSyntaxOff
\def\tifontcase{0}
\IfFileExists{Fonts/TI-84 Plus CE Key Fonts/TI84PlusCEKeys.ttf}
  {\def\tifontcase{1}}{}
\IfFileExists{../Fonts/TI-84 Plus CE Key Fonts/TI84PlusCEKeys.ttf}
  {\def\tifontcase{2}}{}
\ifcase\tifontcase\relax
  % 0: font not found, use boxed fallback keys
  \newcommand{\ti}[1]{\tifallback{#1}}
\or
  % 1: original location (as in preamble.tex)
  \newfontfamily\tifont[
    Path = {Fonts/TI-84 Plus CE Key Fonts/},
    UprightFont = {TI84PlusCEKeys.ttf}
  ]{}
  \newcommand{\ti}[1]{{\tifont #1}}
\or
  % 2: one level up (when compiling inside draft-ch1/)
  \newfontfamily\tifont[
    Path = {../Fonts/TI-84 Plus CE Key Fonts/},
    UprightFont = {TI84PlusCEKeys.ttf}
  ]{}
  \newcommand{\ti}[1]{{\tifont #1}}
\fi

% [FONT SETUP copied from h2-math-notes/aoi/notes/preamble_build.tex (DI notes)]
% [AoI items 17/18/37] GC DISPLAY FONT. Every GC key, command, menu item and
% screen line goes through \gcfont, so the font can be swapped HERE in one place.
%   \gc{fnInt(}      command / menu item / stored variable in the display font
%   \gckey{2nd}      a calculator key (display font in a small key outline)
%   \gcscreen{...}   a calculator screen (display font on a grey-green LCD panel)
% The font file lives in ../Fonts (shared by notes, slides and worksheet).
\def\gcfontfile{VT323-Regular.ttf}
\IfFileExists{Fonts/\gcfontfile}{%
  \newfontfamily\gcfont[Path=Fonts/, Scale=1.32]{\gcfontfile}}{%
  \newcommand\gcfont{\ttfamily}}
\newcommand{\gc}[1]{\mbox{\gcfont #1}}
\newcommand{\gckey}[1]{\tcbox[on line, enhanced, boxsep=0pt, left=2pt, right=2pt,
  top=0.6pt, bottom=0.2pt, boxrule=0.5pt, arc=1.5pt, colback=black!4,
  colframe=black!65]{\gcfont\strut #1}}
\definecolor{lcd}{RGB}{214,222,204}
\newcommand{\gcscreen}[1]{\tcbox[on line, enhanced, boxsep=1pt, left=4pt, right=4pt,
  top=2pt, bottom=2pt, boxrule=1.2pt, arc=1pt, colback=lcd, colframe=black!70]{%
  \gcfont\begin{tabular}{@{}l@{}}#1\end{tabular}}}
% integral template as shown on the TI-84 Plus CE (MathPrint): \gcint{0}{2}{e^(-X²)}{X}
\newcommand{\gcint}[4]{{\gcfont\raisebox{-0.35ex}{\Large ∫}\hspace{-0.15em}%
  \raisebox{-0.2ex}{\begin{tabular}{@{}l@{}}\footnotesize #2\\[-1.3ex]\footnotesize #1\end{tabular}}%
  \hspace{0.1em}(#3)d#4}}
% old key macro kept as an alias so nothing breaks
\renewcommand{\ti}[1]{\gckey{#1}}

% SI units
\usepackage{siunitx}

% setspace
%\usepackage{setspace}
%\onehalfspacing

% For hyperlinks within and out of document
\usepackage{hyperref}
\hypersetup{
    colorlinks = true,
    linkcolor = black,
    urlcolor = blue,
}   

% For captions for pics
\usepackage{capt-of}


% tcolorbox custom environments
% Definition ========================================================
\newtcolorbox[auto counter,number within=section]{definition}[2][drop shadow]{%
enhanced, before skip=5mm, after skip=6mm, % [V3 ADDITION] spacing
colback=blue!5!white,colframe=blue!50!black,fonttitle=\bfseries,
title= Definition~\thetcbcounter. #2,#1}

% Theorem ==========================================================
\newtcolorbox[auto counter,number within=section]{theorem}[2][drop shadow]{%
enhanced, before skip=5mm, after skip=6mm, % [V3 ADDITION] spacing
colback=black!5!white,colframe=black,fonttitle=\bfseries,
title= Theorem~\thetcbcounter. #2,#1}

% Caution box ======================================================
\newtcolorbox{caution}[1][]{enhanced,
before skip=5mm,after skip=6mm, % [V3 ADDITION] was 2mm/3mm
boxrule=1pt,left=11mm,right=2mm,top=3mm,bottom=3mm,
colback=yellow!50,
colframe=yellow!20!black,
sharp corners,rounded corners=southeast,arc is angular,arc=3mm,
underlay={%
\path[fill=tcbcolback!80!black] ([yshift=3mm]interior.south east)--++(-0.4,-0.1)--++(0.1,-0.2);
\path[draw=tcbcolframe,shorten <=-0.05mm,shorten >=-0.05mm] ([yshift=3mm]interior.south east)--++(-0.4,-0.1)--++(0.1,-0.2);
\path[fill=yellow!50!black,draw=none] (interior.south west) rectangle node[white]{\Huge\bfseries !} ([xshift=9mm]interior.north west);
},
drop fuzzy shadow,#1}

% Question box ======================================================
\newtcolorbox{question}[1][]{enhanced,
before skip=5mm,after skip=6mm, % [V3 ADDITION] was 2mm/3mm
boxrule=1pt,left=11mm,right=2mm,top=3mm,bottom=3mm,
colback=blue!20!white,
colframe=blue!20!black,
sharp corners,rounded corners=southeast,arc is angular,arc=3mm,
underlay={%
\path[fill=tcbcolback!80!black] ([yshift=3mm]interior.south east)--++(-0.4,-0.1)--++(0.1,-0.2);
\path[draw=tcbcolframe,shorten <=-0.05mm,shorten >=-0.05mm] ([yshift=3mm]interior.south east)--++(-0.4,-0.1)--++(0.1,-0.2);
\path[fill=blue!50!black,draw=none] (interior.south west) rectangle node[white]{\Huge\bfseries ?} ([xshift=9mm]interior.north west);
},
drop fuzzy shadow,#1}

% Information box ======================================================
\newtcolorbox{info}[1][]{enhanced,
before skip=5mm,after skip=6mm, % [V3 ADDITION] was 2mm/3mm
boxrule=1pt,left=11mm,right=2mm,top=3mm,bottom=3mm,
colback=cyan!20!white,
colframe=cyan!20!black,
sharp corners,rounded corners=southeast,arc is angular,arc=3mm,
underlay={%
\path[fill=tcbcolback!80!black] ([yshift=3mm]interior.south east)--++(-0.4,-0.1)--++(0.1,-0.2);
\path[draw=tcbcolframe,shorten <=-0.05mm,shorten >=-0.05mm] ([yshift=3mm]interior.south east)--++(-0.4,-0.1)--++(0.1,-0.2);
\path[fill=cyan!50!black,draw=none] (interior.south west) rectangle node[white]{\Huge\bfseries i} ([xshift=9mm]interior.north west);
},
drop fuzzy shadow,#1}

% Mistake box ======================================================
\newtcolorbox{mistake}[1][]{enhanced,
before skip=5mm,after skip=6mm, % [V3 ADDITION] was 2mm/3mm
boxrule=1pt,left=11mm,right=2mm,top=3mm,bottom=3mm,
colback=red!20!white,
colframe=red!20!black,
sharp corners,rounded corners=southeast,arc is angular,arc=3mm,
underlay={%
\path[fill=tcbcolback!80!black] ([yshift=3mm]interior.south east)--++(-0.4,-0.1)--++(0.1,-0.2);
\path[draw=tcbcolframe,shorten <=-0.05mm,shorten >=-0.05mm] ([yshift=3mm]interior.south east)--++(-0.4,-0.1)--++(0.1,-0.2);
\path[fill=red!50!black,draw=none] (interior.south west) rectangle node[white]{\Huge\bfseries $\times$} ([xshift=9mm]interior.north west);
},
drop fuzzy shadow,#1}



% [V3 ADDITION] Worked example box ==================================
% Own counter, numbered within the section (Example 11.2.1, 11.2.2, ...).
% For per-subsection numbering (11.2.1.1) change `section' to `subsection'. Blue `question' boxes are kept for thought-provoking questions
% only; worked examples always use this box. Usage:
%   \begin{example} ...question text... \end{example}
%   \solution{...}
\newtcolorbox[auto counter,number within=section]{example}[1][]{%
enhanced, breakable=false, before skip=5mm, after skip=4mm,
colback=teal!5!white, colframe=teal!55!black, boxrule=0.8pt,
left=3mm, right=3mm, top=2mm, bottom=2mm,
sharp corners, rounded corners=northwest, arc=3mm,
fonttitle=\bfseries, coltitle=white,
title=Example~\thetcbcounter, attach boxed title to top left={yshift=-2mm, xshift=3mm},
boxed title style={colback=teal!55!black, sharp corners, rounded corners=northwest, arc=2mm},
drop fuzzy shadow, #1}

% [V3 ADDITION] Spacing between sections, subsections, parts, paragraphs =
% (fncychap's chapter heading is untouched; only \section and below.)
\makeatletter
\renewcommand\section{\@startsection{section}{1}{\z@}%
  {-2ex \@plus -.5ex \@minus -.2ex}{3.2ex \@plus .3ex}%
  {\normalfont\LARGE\bfseries\color{blue!50!black}}}
\renewcommand\subsection{\@startsection{subsection}{2}{\z@}%
  {-1.5ex \@plus -.5ex \@minus -.2ex}{2.6ex \@plus .3ex}%
  {\normalfont\Large\bfseries}}
\renewcommand\subsubsection{\@startsection{subsubsection}{3}{\z@}%
  {-4ex \@plus -1ex \@minus -.5ex}{1.6ex \@plus .2ex}%
  {\normalfont\large\bfseries\color{teal!50!black}}}
\makeatother
\setlength{\parindent}{0pt}
\setlength{\parskip}{0.9ex plus 0.3ex minus 0.2ex}

% Headers and Footers
\usepackage{fancyhdr,lastpage}
% [V3 R2] grey left-margin line on the out-of-syllabus pages (item 13)
\usepackage{eso-pic}
\newcommand{\oosmargin}{\AddToShipoutPictureBG{\put(12mm,0){\color{gray!75}\rule{1.6pt}{\paperheight}}}}
\pagestyle{fancy}
\lhead{Maximus Pung Jun}
% [BUILD COPY CHANGE 2] Header bug fix: \chaptername is the word "Chapter",
% so the original header printed "Chapter 1: Chapter". \chaptermark now stores
% the bare chapter title in \leftmark, and the header uses \leftmark.
\renewcommand{\chaptermark}[1]{\markboth{#1}{}}
\rhead{Chapter \thechapter: \leftmark}
% \cfoot{Page \thepage\  of \pageref{LastPage}} % buggy, does not appear on chapter pages.

% [V3 R2] tick and cross markers for mistake boxes (G3)
\newcommand{\wrong}{\textcolor{red!70!black}{\ensuremath{\boldsymbol{\times}}}}
\newcommand{\cmark}{\textcolor{green!50!black}{\ensuremath{\checkmark}}}

\newcommand{\xmark}{\wrong}   % alias: red cross
% [V3 R2] out-of-syllabus helpers (DE notes item 13)
%   \oosmargin    : grey line in the left margin on every following page
%   \oosmarginoff : stop the margin line
%   \oosnote      : grey banner "Out of syllabus. For your interest only."
%   outofsyllabus : environment = new page + margin line + banner; ends with
%                   \clearpage and switches the margin line off again
\newcommand{\oosmarginoff}{\ClearShipoutPictureBG}
\newcommand{\oosnote}[1][Nothing in this section is examined in H2 Mathematics.]{\par\noindent\colorbox{gray!15}{\parbox{\dimexpr\textwidth-2\fboxsep}{\textbf{Out of syllabus. For your interest only.} #1}}\par}
\newenvironment{outofsyllabus}{\clearpage\oosmargin}{\clearpage\oosmarginoff}

% Notation as in Maximus's notes: italic d, display style everywhere (G1).
% Fractions inside exponents stay small, so never use these macros inside ^{}.
\newcommand{\dydx}{\dfrac{dy}{dx}}
\newcommand{\dvdx}{\dfrac{dv}{dx}}
\newcommand{\dudx}{\dfrac{du}{dx}}
\newcommand{\ddydx}{\dfrac{d^{2}y}{dx^{2}}}
\newcommand{\ddt}[1]{\dfrac{d#1}{dt}}
% Small syllabus flag (no longer used in the DE notes, kept for other chapters)
\newcommand{\syllabusflag}[1]{\par\noindent{\footnotesize\color{red!70!black}\textbf{Syllabus note.} #1}\par}


% Teacher's copy toggle ================================================
% [BUILD COPY CHANGE 3] \ifx compares the whole word (the original
% \if\conform test was always false). DE_notes_student.tex predefines
% \conform as HideSolutions.
% A student copy defines \conform as HideSolutions BEFORE loading this preamble.
\providecommand\conform{ShowSolutions}
\newcommand\ShowSolutionsKey{ShowSolutions}
% [V3 ADDITION] In the student copy the hidden solution leaves stretchable
% blank space (\vfill), so the free space on each page is shared out as
% working space below the examples instead of collecting at the bottom.
\newcommand\solution[1]{\ifx\conform\ShowSolutionsKey
  \par\noindent\textbf{Solution.} #1\par\medskip\else\par\vspace{\workspace}\fi}
% [V3 R2] minimum working space left under each example in the student copy
\newlength\workspace\setlength\workspace{3cm plus 1fill}


\newcommand{\NN}{\mathbb{N}}
\newcommand{\ZZ}{\mathbb{Z}}
\newcommand{\QQ}{\mathbb{Q}}
\newcommand{\RR}{\mathbb{R}}
\newcommand{\CC}{\mathbb{C}}

